% =========================================================
% Algebras of Sets - Rings of Sets
% =========================================================
\section{Rings of Sets}

\begin{toolkitbox}{Toolkit --- Boolean Ring Examples}
\textbf{Concept family:} Concrete Boolean ring examples and structural
constraints.

\medskip
\textbf{Atomic items in this section:}
\begin{itemize}
  \item Remark: No Boolean ring of order 3 exists
  \item Example: The four-element Boolean ring $\mathbf{2}^2$
  \item Remark: Boolean rings have order $2^n$
\end{itemize}

\medskip
\textbf{Key constraint:} $(R,+)$ in any Boolean ring is an elementary
abelian 2-group, so $|R| = 2^n$ for some $n \ge 0$.
\end{toolkitbox}


\begin{toolkitbox}{Toolkit --- Boolean Rings}
\textbf{Concept family:} Boolean rings (Halmos--Givant §1) and their
immediate consequences from the single idempotence axiom.

\medskip
\textbf{Atomic items in this section:}
\begin{itemize}
  \item Definition: Boolean Ring
  \item Definition: Idempotent Element
  \item Definition: Characteristic 2
  \item Lemma: Self Additive Inverse ($p+p=0$)
  \item Lemma: Self Negation ($-p=p$)
\end{itemize}

\medskip
\textbf{Key predicate:}
\[
  \operatorname{BooleanRing}(R,+,\cdot,-,0,1)
  \;\colonequiv\;
  \operatorname{Ring}(R,+,\cdot,-,0,1)
  \;\land\;
  \operatorname{IdempotentRing}(R,\cdot).
\]
\textbf{Central theorem chain:} Idempotence $\Rightarrow$ Characteristic 2
$\Rightarrow$ Self Negation ($-p=p$) $\Rightarrow$ Commutativity of $\cdot$.
\end{toolkitbox}


\begin{toolkitbox}{Toolkit --- Rings of Sets}
\textbf{Concept family:} The ring axioms (Halmos--Givant §1) as the algebraic
foundation for Boolean algebras of sets.

\medskip
\textbf{Atomic items in this section:}
\begin{itemize}
  \item Definition: Ring
\end{itemize}

\medskip
\textbf{Key predicate:}
\[
  \operatorname{Ring}(R,+,\cdot,-,0,1)
  \;\colonequiv\;
  \operatorname{AbelianGroup}(R,+,-,0)
  \;\land\;
  \operatorname{Monoid}(R,\cdot,1)
  \;\land\;
  \operatorname{Distributes}(\cdot,+,R).
\]
\textbf{Axiom numbering} follows Halmos--Givant; axiom (4) (commutativity of
multiplication) is not a ring axiom but a theorem in the Boolean setting.
\end{toolkitbox}


% =========================================================
% Definition: Ring
% Section: Rings of Sets — Volume I
% Source: Halmos & Givant, Introduction to Boolean Algebras, §1
% =========================================================

\begin{tcolorbox}[colback=gray!6, colframe=gray!40, arc=2pt,
  left=6pt, right=6pt, top=4pt, bottom=4pt,
  title={\small\textbf{Toolkit --- Rings of Sets}},
  fonttitle=\small\bfseries]
\textbf{Concept family:} The ring axioms (Halmos--Givant §1) as the algebraic
foundation for Boolean algebras of sets.

\medskip
\textbf{Atomic items in this section:}
\begin{itemize}
  \item Definition: Ring
\end{itemize}

\medskip
\textbf{Key predicate:}
\[
  \operatorname{Ring}(R,+,\cdot,-,0,1)
  \;\colonequiv\;
  \operatorname{AbelianGroup}(R,+,-,0)
  \;\land\;
  \operatorname{Monoid}(R,\cdot,1)
  \;\land\;
  \operatorname{Distributes}(\cdot,+,R).
\]
\textbf{Axiom numbering} follows Halmos--Givant; axiom (4) (commutativity of
multiplication) is not a ring axiom but a theorem in the Boolean setting.
\end{tcolorbox}

\begin{tcolorbox}[colback=defbox, colframe=defborder, arc=2pt,
  left=6pt, right=6pt, top=4pt, bottom=4pt,
  title={\small\textbf{Definition --- Ring}},
  fonttitle=\small\bfseries]
\begin{definition}[Ring]
\label{def:ring}
A \textbf{ring} is a non-empty set $R$ together with two binary
operations $+$ and $\cdot$ on $R$, a unary operation $-$ on $R$,
and two distinguished elements $0, 1 \in R$, satisfying the
following axioms for all $p, q, r \in R$:
\begin{enumerate}
  \item[(1)] $p + (q + r) = (p + q) + r$
             \hfill\textit{(additive associativity)}
  \item[(2)] $p \cdot (q \cdot r) = (p \cdot q) \cdot r$
             \hfill\textit{(multiplicative associativity)}
  \item[(3)] $p + q = q + p$
             \hfill\textit{(additive commutativity)}
  \item[(5)] $p + 0 = p$
             \hfill\textit{(additive identity)}
  \item[(6)] $p \cdot 1 = p$
             \hfill\textit{(multiplicative identity)}
  \item[(7)] $p + (-p) = 0$
             \hfill\textit{(additive inverse)}
  \item[(8)] $p \cdot (q + r) = p \cdot q + p \cdot r$
             \hfill\textit{(left distributivity)}
  \item[(9)] $(q + r) \cdot p = q \cdot p + r \cdot p$
             \hfill\textit{(right distributivity)}
\end{enumerate}
Axiom numbering follows Halmos--Givant~\cite{HalmosGivant2009};
axiom~(4) is commutativity of multiplication, which is not a ring
axiom but a theorem in the Boolean setting.
A ring possessing a distinguished element $1$ satisfying
axiom~(6) is called a \textbf{ring with unit}.
A ring satisfying $p \cdot q = q \cdot p$ for all $p, q \in R$
is called a \textbf{commutative ring}.
\end{definition}
\end{tcolorbox}

\begin{remark*}[Standard quantified statement]
\[
\begin{aligned}
(R,+,\cdot,-,0,1)\ \text{is a ring}
\;\Longleftrightarrow\;
  &\forall p,q,r \in R{:} \\
  &\quad p + (q + r) = (p + q) + r \\
  &\;\land\quad p \cdot (q \cdot r) = (p \cdot q) \cdot r \\
  &\;\land\quad p + q = q + p \\
  &\;\land\quad p + 0 = p \\
  &\;\land\quad p \cdot 1 = p \\
  &\;\land\quad p + (-p) = 0 \\
  &\;\land\quad p \cdot (q + r) = p \cdot q + p \cdot r \\
  &\;\land\quad (q + r) \cdot p = q \cdot p + r \cdot p.
\end{aligned}
\]
\end{remark*}

\begin{remark*}[Definition predicate reading]
\[
\operatorname{Ring}(R,+,\cdot,-,0,1)
\;\Longleftrightarrow\;
\operatorname{AbelianGroup}(R,+,-,0)
\;\land\;
\operatorname{Monoid}(R,\cdot,1)
\;\land\;
\operatorname{Distributes}(\cdot,+,R)
\]
\end{remark*}

\begin{remark*}[Negated quantified statement]
\[
\begin{aligned}
(R,+,\cdot,-,0,1)\ \text{is \textbf{not} a ring}
\;\Longleftrightarrow\;
  &\exists p,q,r \in R{:} \\
  &\quad p + (q + r) \neq (p + q) + r \\
  &\;\lor\quad p \cdot (q \cdot r) \neq (p \cdot q) \cdot r \\
  &\;\lor\quad p + q \neq q + p \\
  &\;\lor\quad p + 0 \neq p \\
  &\;\lor\quad p \cdot 1 \neq p \\
  &\;\lor\quad p + (-p) \neq 0 \\
  &\;\lor\quad p \cdot (q + r) \neq p \cdot q + p \cdot r \\
  &\;\lor\quad (q + r) \cdot p \neq q \cdot p + r \cdot p.
\end{aligned}
\]
\end{remark*}

\begin{remark*}[Failure mode decomposition]
The eight axioms partition into three structural groups:
\[
\underbrace{(1),(3),(5),(7)}_{\displaystyle\operatorname{AbelianGroup}(R,+,-,0)}
\qquad
\underbrace{(2),(6)}_{\displaystyle\operatorname{Monoid}(R,\cdot,1)}
\qquad
\underbrace{(8),(9)}_{\displaystyle\operatorname{Distributes}(\cdot,+,R)}
\]
A structure violating any axiom in the first group lacks a coherent
additive group.
A structure violating any axiom in the second group lacks a
multiplicative identity or associativity.
A structure violating the third group possesses two individually
well-behaved operations that do not interact coherently;
distributivity is the sole structural bridge between $+$ and
$\cdot$, and its failure is sufficient to destroy the ring property.
\end{remark*}

\begin{remark*}[Interpretation]
A ring abstracts the arithmetic of the integers.
The additive structure is always a commutative group.
The multiplicative structure is a monoid --- associative with
unit --- but need not be commutative; axiom~(4) of
Halmos--Givant, commutativity of multiplication, is absent from
the ring axioms.
In the Boolean setting, multiplicative commutativity is not
assumed but derived as a theorem from idempotence alone.
Distributivity encodes the coherent interaction between the
two operations: it is the axiom that makes $\cdot$ and $+$
form a single algebraic system rather than two independent ones.
\end{remark*}

\NoLocalDependencies

% ---------------------------------------------------------
% Section: Boolean Rings
% ---------------------------------------------------------

% ---------------------------------------------------------
% Definition: Boolean Ring
% Source: Halmos & Givant, Introduction to Boolean Algebras, \S{}1
% ---------------------------------------------------------

\begin{definitionbox}{Definition --- Boolean Ring}
\begin{definition}[Boolean Ring]
\label{def:boolean-ring}
A \textbf{Boolean ring} is an idempotent ring with unit.

That is, a Boolean ring is a ring $(R, +, \cdot, -, 0, 1)$
satisfying the additional axiom:
\[
p \cdot p = p \qquad \text{for all } p \in R. \tag{11}
\]
An element $p$ satisfying $p \cdot p = p$ is called
\textbf{idempotent}.
A ring in which every element is idempotent is called an
\textbf{idempotent ring}.

The official axiom set for a Boolean ring consists of
axioms~(1)--(3), (5), (6), (8)--(11), where axiom~(7)
$($the additive inverse law$)$ is replaced by
axiom~(10) $($characteristic~$2)$, which is derived as a theorem.
Axiom numbering follows Halmos--Givant~\cite{HalmosGivant2009}.
\end{definition}
\end{definitionbox}

\begin{remark*}[Standard quantified statement]
\[
\begin{aligned}
\operatorname{BooleanRing}(R,+,\cdot,-,0,1)
&\;\Longleftrightarrow\;
  \operatorname{Ring}(R,+,\cdot,-,0,1)
  \;\land\;
  \forall p \in R{:}\; p \cdot p = p.
\end{aligned}
\]
\end{remark*}

\begin{remark*}[Definition predicate reading]
\[
\operatorname{BooleanRing}(R,+,\cdot,-,0,1)
\;\Longleftrightarrow\;
\operatorname{Ring}(R,+,\cdot,-,0,1)
\;\land\;
\operatorname{IdempotentRing}(R,\cdot)
\]
\end{remark*}

\begin{remark*}[Negated quantified statement]
\[
\neg\,\operatorname{BooleanRing}(R,+,\cdot,-,0,1)
\;\Longleftrightarrow\;
\neg\,\operatorname{Ring}(R,+,\cdot,-,0,1)
\;\lor\;
\exists p \in R{:}\; p \cdot p \neq p.
\]
\end{remark*}

\begin{remark*}[Failure modes]
The failure modes are exactly the ways that one of the defining structural
conditions can fail.
\end{remark*}

\begin{remark*}[Failure mode decomposition]
A structure fails to be a Boolean ring in exactly one of two ways:
\[
\underbrace{\neg\,\operatorname{Ring}(R,+,\cdot,-,0,1)}_{\text{fails a ring axiom}}
\qquad\lor\qquad
\underbrace{\exists p \in R{:}\; p \cdot p \neq p}_{\text{some element is not idempotent}}
\]
The second failure mode is the structurally interesting one.
A ring in which even one element fails $p \cdot p = p$ is not Boolean.
The simplest example: the ring $\mathbb{Z}$ of integers fails
idempotence at every element except $0$ and $1$.
\end{remark*}

\begin{remark*}[Canonical examples]
\hfill
\begin{itemize}
  \item $\mathbf{2} = \{0,1\}$ with addition mod~$2$ and ordinary
        multiplication is the simplest non-degenerate Boolean ring.
  \item $\mathbf{2}^n$, the set of all $n$-tuples of elements of
        $\mathbf{2}$, with coordinatewise addition and multiplication,
        is a Boolean ring with $2^n$ elements.
        The zero is $(0,\ldots,0)$ and the unit is $(1,\ldots,1)$.
  \item $(\mathcal{P}(X), \triangle, \cap)$, the power set of any
        set $X$ with symmetric difference as addition and intersection
        as multiplication, is a Boolean ring.
        Under the correspondence $A \leftrightarrow \mathbf{1}_A$
        (the characteristic function of $A$), this ring is isomorphic
        to $\mathbf{2}^X$.
\end{itemize}
\end{remark*}

\begin{remark*}[Interpretation]
A Boolean ring is the minimal algebraic structure that forces
the full Boolean calculus to emerge from a single additional
axiom --- idempotence of multiplication.
The ring axioms supply the additive abelian group and the
multiplicative monoid with distributivity.
Idempotence then forces characteristic~$2$ (every element
equals its own additive inverse), forces commutativity of
multiplication, and identifies the ring with a field of sets
via the representation theorem (Halmos--Givant, Chapter~22).
The abstract definition thus captures exactly the algebra of
set-theoretic operations: symmetric difference plays the role
of addition and intersection plays the role of multiplication.
\end{remark*}

\begin{dependencies}
\begin{itemize}
  \item \hyperref[def:algebraic-ring]{Definition: Ring}
\end{itemize}
\end{dependencies}


% =========================================================
% Section: Boolean Rings --- Definitions and First Properties
% Chapter: Boolean Rings and Boolean Algebras --- Volume I
% Source: Halmos & Givant, Introduction to Boolean Algebras, \S{}1
% =========================================================

% ---------------------------------------------------------
% The Two-Element Boolean Ring 2
% ---------------------------------------------------------

\begin{example*}[The two-element Boolean ring]
The simplest non-degenerate Boolean ring is the two-element ring
$\mathbf{2} = \{0, 1\}$.
Its addition and multiplication are given by the following tables.

\begin{center}
\begin{tabular}{c|cc}
$+$ & $0$ & $1$ \\ \hline
$0$ & $0$ & $1$ \\
$1$ & $1$ & $0$
\end{tabular}
\qquad\qquad
\begin{tabular}{c|cc}
$\cdot$ & $0$ & $1$ \\ \hline
$0$ & $0$ & $0$ \\
$1$ & $0$ & $1$
\end{tabular}
\end{center}

\noindent
Addition in $\mathbf{2}$ is addition modulo $2$: $1 + 1 = 0$.
Multiplication in $\mathbf{2}$ is ordinary multiplication of
integers restricted to $\{0, 1\}$.
Every element of $\mathbf{2}$ satisfies $p + p = 0$ and
$p \cdot p = p$, making $\mathbf{2}$ the canonical example against
which both axioms~(10) and~(11) are verified.
\end{example*}

% ---------------------------------------------------------
% Definition: Idempotent Element
% ---------------------------------------------------------

\begin{definitionbox}{Definition --- Idempotent Element}
\begin{definition}[Idempotent Element]
\label{def:idempotent-element}
An element $p$ of a ring $R$ is called \textbf{idempotent} if
\[
p \cdot p = p.
\]
A ring in which every element is idempotent is called an
\textbf{idempotent ring}.
\end{definition}
\end{definitionbox}

\begin{remark*}[Standard quantified statement]
\[
\begin{aligned}
\operatorname{IdempotentElement}(p, R, \cdot)
&\;\Longleftrightarrow\; p \cdot p = p. \\[4pt]
\operatorname{IdempotentRing}(R, \cdot)
&\;\Longleftrightarrow\; \forall p \in R{:}\; p \cdot p = p.
\end{aligned}
\]
\end{remark*}

\begin{remark*}[Definition predicate reading]
\[
\operatorname{IdempotentElement}(p,R,\cdot)
\qquad
\operatorname{IdempotentRing}(R,\cdot)
\;\Longleftrightarrow\;
\forall p \in R{:}\;
\operatorname{IdempotentElement}(p,R,\cdot)
\]
\end{remark*}

\begin{remark*}[Negated quantified statement]
\[
\neg\,\operatorname{IdempotentElement}(p,R,\cdot)
\;\Longleftrightarrow\;
p \cdot p \neq p.
\]
\end{remark*}

\begin{remark*}[Interpretation]
Idempotence is the single additional axiom that forces Boolean
structure.  In the two-element ring $\mathbf{2}$, verification is
immediate: $0 \cdot 0 = 0$ and $1 \cdot 1 = 1$.  In a power set
algebra $\mathcal{P}(X)$, idempotence of intersection --- $A \cap A
= A$ --- is the set-theoretic counterpart.  The force of the axiom
is not in these examples but in what it implies for an arbitrary
ring: characteristic~2 and commutativity of multiplication both
follow from idempotence alone, as the theorems below establish.
\end{remark*}

\begin{dependencies}
\begin{itemize}
  \item \hyperref[def:algebraic-ring]{Definition: Ring}
\end{itemize}
\end{dependencies}

% ---------------------------------------------------------
% Definition: Characteristic 2
% ---------------------------------------------------------

\begin{definitionbox}{Definition --- Characteristic 2}
\begin{definition}[Characteristic 2]
\label{def:characteristic-2}
A ring $R$ is said to have \textbf{characteristic~$2$} if
\[
p + p = 0 \qquad \text{for all } p \in R. \tag{10}
\]
\end{definition}
\end{definitionbox}

\begin{remark*}[Standard quantified statement]
\[
\operatorname{Characteristic2}(R,+,0)
\;\Longleftrightarrow\;
\forall p \in R{:}\; p + p = 0.
\]
\end{remark*}

\begin{remark*}[Definition predicate reading]
\[
\operatorname{Characteristic2}(R,+,0)
\;\Longleftrightarrow\;
\forall p \in R{:}\;
\operatorname{IdempotentElement}(p, R, +)
\]
\end{remark*}

\begin{remark*}[Negated quantified statement]
\[
\neg\,\operatorname{Characteristic2}(R,+,0)
\;\Longleftrightarrow\;
\exists p \in R{:}\; p + p \neq 0.
\]
\end{remark*}

\begin{remark*}[Interpretation]
Characteristic~2 means that every element is its own additive
inverse: $p + p = 0$ says precisely $p = -p$.  The operation of
negation therefore becomes the identity map --- applying $-$ to
any element returns that element unchanged.  Halmos observes that
in a ring of characteristic~2, the minus sign for additive inverses
is never needed and is henceforth dropped.  Axiom~(7),
$p + (-p) = 0$, is replaced in the official Boolean ring axiom set
by the stronger equation~(10), $p + p = 0$, which is a theorem
derived from idempotence.
\end{remark*}

\begin{dependencies}
\begin{itemize}
  \item \hyperref[def:boolean-ring]{Definition: Boolean Ring}
  \item \hyperref[def:idempotent-element]{Definition: Idempotent Element}
\end{itemize}
\end{dependencies}

% ---------------------------------------------------------
% Lemma: Self Additive Inverse (Characteristic 2)
% ---------------------------------------------------------

\begin{lemmabox}{Lemma}
\begin{lemma}[Self Additive Inverse]
\label{lem:self-additive-inverse}
In any Boolean ring $R$,
\[
p + p = 0 \qquad \text{for all } p \in R.
\]

\smallskip
\noindent
\hyperref[prf:self-additive-inverse]{\textit{Go to proof.}}
\end{lemma}
\end{lemmabox}

\begin{remark*}[Standard quantified statement]
\[
\forall p \in R{:}\; p + p = 0.
\]
\end{remark*}

\begin{remark*}[Definition predicate reading]
\[
\operatorname{BooleanRing}(R,+,\cdot,-,0,1)
\;\Longrightarrow\;
\operatorname{Characteristic2}(R,+,0).
\]
\end{remark*}

\begin{remark*}[Interpretation]
This is Halmos consequence~(a): idempotence of multiplication
forces the additive structure to have characteristic~2.  The proof
expands $(p+q)^2$ by distributivity and idempotence, cancels $p$
and $q$, then sets $p = q$ to obtain $p + p = 0$.  The Lean~4
formalisation is \texttt{AdditiveIdempotence} in
\texttt{BooleanAlgebraDefinitions.lean}.
\end{remark*}

\begin{dependencies}
\begin{itemize}
  \item \hyperref[def:boolean-ring]{Definition: Boolean Ring}
  \item \hyperref[def:characteristic-2]{Definition: Characteristic 2}
\end{itemize}
\end{dependencies}

% ---------------------------------------------------------
% Lemma: Self Negation
% ---------------------------------------------------------

\begin{lemmabox}{Lemma}
\begin{lemma}[Self Negation]
\label{lem:self-negation}
In any Boolean ring $R$,
\[
-p = p \qquad \text{for all } p \in R.
\]

\smallskip
\noindent
\hyperref[prf:self-negation]{\textit{Go to proof.}}
\end{lemma}
\end{lemmabox}

\begin{remark*}[Standard quantified statement]
\[
\forall p \in R{:}\; -p = p.
\]
\end{remark*}

\begin{remark*}[Interpretation]
This follows immediately from \ref{lem:self-additive-inverse}.
Since $p + p = 0$ and $p + (-p) = 0$ both hold, the uniqueness of
additive inverses gives $-p = p$.  The negation operation is
therefore the identity map on any Boolean ring.  Halmos uses this
to drop the minus sign entirely from his notation for the rest of
the book: since $-p = p$, writing $-p$ and writing $p$ are the
same thing.
\end{remark*}

\begin{dependencies}
\begin{itemize}
  \item \hyperref[lem:self-additive-inverse]{Lemma: Self Additive Inverse}
\end{itemize}
\end{dependencies}

% =========================================================
% Scratch: No 3-Element Boolean Ring — Example Note
% Chapter: Boolean Rings and Boolean Algebras — Volume I
% Source: Halmos & Givant, Introduction to Boolean Algebras, §1
%
% STATUS: scratch file — move into chapter structure when
%         volume-i/boolean-rings/ directory is created.
%
% Depends on:
%   \ref{lem:self-additive-inverse}  (p + p = 0)
%   colors: defbox/defborder, propbox/propborder (common/colors.tex)
%   environments: example*, remark* (common/environments.tex)
%   packages: booktabs (for \toprule etc.)
% =========================================================

% ---------------------------------------------------------
% Remark: No Boolean Ring of Order 3 Exists
% ---------------------------------------------------------

\begin{remark*}[No Boolean ring of order 3]
A Boolean ring of order~3 does not exist.
The additive group $(R,+)$ of any Boolean ring satisfies
$p + p = 0$ for all $p$, by Lemma~\ref{lem:self-additive-inverse}.
Every element therefore has additive order dividing~2, so
$(R,+)$ is an elementary abelian $2$-group.
Elementary abelian $2$-groups have order $2^n$ for some $n \ge 0$.
Since $3$ is not a power of~$2$, no Boolean ring of order~$3$
can exist.
\end{remark*}

% ---------------------------------------------------------
% Example: The Four-Element Boolean Ring 2^2
% ---------------------------------------------------------

\begin{example*}[The four-element Boolean ring $\mathbf{2}^2$]
\label{ex:boolean-ring-four-elements}
The smallest Boolean ring with more than two elements is
$\mathbf{2}^2 = \{(0,0),\,(0,1),\,(1,0),\,(1,1)\}$,
the set of ordered pairs from $\mathbf{2} = \{0,1\}$.
Addition and multiplication are defined coordinatewise:
\[
  (p_0,p_1) + (q_0,q_1) \;=\; (p_0+q_0,\;p_1+q_1),
  \qquad
  (p_0,p_1) \cdot (q_0,q_1) \;=\; (p_0 \cdot q_0,\;p_1 \cdot q_1),
\]
where each coordinate is computed in $\mathbf{2}$.
The zero element is $(0,0)$ and the unit is $(1,1)$.

\medskip
\noindent\textbf{Addition table.}
\begin{center}
\renewcommand{\arraystretch}{1.25}
\begin{tabular}{c|cccc}
$+$     & $(0,0)$ & $(0,1)$ & $(1,0)$ & $(1,1)$ \\ \hline
$(0,0)$ & $(0,0)$ & $(0,1)$ & $(1,0)$ & $(1,1)$ \\
$(0,1)$ & $(0,1)$ & $(0,0)$ & $(1,1)$ & $(1,0)$ \\
$(1,0)$ & $(1,0)$ & $(1,1)$ & $(0,0)$ & $(0,1)$ \\
$(1,1)$ & $(1,1)$ & $(1,0)$ & $(0,1)$ & $(0,0)$
\end{tabular}
\end{center}

\medskip
\noindent\textbf{Multiplication table.}
\begin{center}
\renewcommand{\arraystretch}{1.25}
\begin{tabular}{c|cccc}
$\cdot$ & $(0,0)$ & $(0,1)$ & $(1,0)$ & $(1,1)$ \\ \hline
$(0,0)$ & $(0,0)$ & $(0,0)$ & $(0,0)$ & $(0,0)$ \\
$(0,1)$ & $(0,0)$ & $(0,1)$ & $(0,0)$ & $(0,1)$ \\
$(1,0)$ & $(0,0)$ & $(0,0)$ & $(1,0)$ & $(1,0)$ \\
$(1,1)$ & $(0,0)$ & $(0,1)$ & $(1,0)$ & $(1,1)$
\end{tabular}
\end{center}

\medskip
\noindent\textbf{Verification of idempotence.}
Every element squares to itself:
\begin{center}
\renewcommand{\arraystretch}{1.2}
\begin{tabular}{cc}
\toprule
$p$     & $p \cdot p$ \\ \midrule
$(0,0)$ & $(0,0)$     \\
$(0,1)$ & $(0,1)$     \\
$(1,0)$ & $(1,0)$     \\
$(1,1)$ & $(1,1)$     \\ \bottomrule
\end{tabular}
\end{center}

\medskip
\noindent\textbf{Verification of characteristic~$2$.}
Every element is its own additive inverse:
\begin{center}
\renewcommand{\arraystretch}{1.2}
\begin{tabular}{cc}
\toprule
$p$     & $p + p$ \\ \midrule
$(0,0)$ & $(0,0)$ \\
$(0,1)$ & $(0,0)$ \\
$(1,0)$ & $(0,0)$ \\
$(1,1)$ & $(0,0)$ \\ \bottomrule
\end{tabular}
\end{center}
\end{example*}

% ---------------------------------------------------------
% Remark: The Pattern — Boolean Rings Have Order 2^n
% ---------------------------------------------------------

\begin{remark*}[Boolean rings have order $2^n$]
Every finite Boolean ring is isomorphic to $\mathbf{2}^n$ for
some $n \ge 0$, and therefore has exactly $2^n$ elements.
This follows from the representation theorem
(Halmos--Givant, Chapter~22): every Boolean algebra is isomorphic
to a field of sets, and every finite field of sets is isomorphic
to $\mathcal{P}(X)$ for some finite set $X$ with $|X| = n$,
giving $|\mathcal{P}(X)| = 2^n$.
The permissible finite sizes are $1, 2, 4, 8, 16, \ldots$
and no other cardinality is achievable.
\end{remark*}

